\chapter{Faits stylisés des rendements}\label{ch:rs:stylised-facts-of-returns}

Le 19 octobre 1987, le Dow Jones Industrial Average a chuté de 22.6~\% en une séance. Tout le marché boursier américain, pondéré par la capitalisation et
en excès des bons du Trésor, a chuté de 17.4~\%~: 22 fois l'écart-type de ses rendements quotidiens des cinq années précédentes. Sous une loi normale,
une telle journée a une probabilité trop faible pour mériter d'être notée~; sur un siècle de données quotidiennes, c'est l'un des 103 jours situés au-delà de cinq
écarts-types, alors qu'une loi normale n'en attendrait que 0.015. Les rendements réels ne sont pas normaux, et les manières dont ils ne le sont pas sont
assez régulières pour avoir un nom. Ce chapitre énonce les \omterm{def:rs:stylised-facts-of-returns:fact}{faits stylisés} des rendements et les mesure sur un siècle de données américaines~; il
construit ensuite le marché actions synthétique du livre et vérifie quels faits il reproduit et lesquels il manque, car tout chapitre ultérieur
qui utilise le simulateur hérite des deux.

\section{Queues épaisses et gaussianité par agrégation}

\begin{definition}[Fait stylisé]\label{def:rs:stylised-facts-of-returns:fact}
Un \emph{fait stylisé}\index{fait stylisé} est une propriété statistique des rendements observée, qualitativement, sur de nombreux actifs,
marchés et périodes, énoncée sans la précision qui la rendrait propre à l'un d'entre eux.
\end{definition}

Cont (2001) en recense onze. Les deux premiers sont l'absence d'autocorrélation linéaire des rendements au-delà d'horizons intrajournaliers très courts,
et les queues épaisses~: les queues de la distribution des rendements décroissent comme une puissance, avec un indice de queue (livre~4, chapitre~15) supérieur à
deux et inférieur à cinq pour la plupart des jeux de données. Le rendement quotidien excédentaire du marché américain de juillet 1926 à juillet 2026 (26\,296 jours, d'après la
Kenneth French data library) présente un kurtosis de 19.1~; l'estimateur de Hill sur les 500 plus fortes pertes donne un indice de queue de 2.90, et
sur les 500 plus forts gains 2.58. Gopikrishnan et ses coauteurs (1998) trouvent un exposant d'environ trois dans quarante millions de variations de cours, transaction par transaction,
d'actions américaines, la «~loi cubique inverse~». La \Cref{fig:rs:stylised-facts-of-returns:hist} montre les rendements quotidiens du siècle
face à une loi normale de même écart-type.

\begin{omfigure}
\begin{tikzpicture}
\begin{semilogyaxis}[width=10.5cm,height=5.4cm,xlabel={rendement quotidien excédentaire en écarts-types},ylabel={fréquence (échelle logarithmique)},
  tick label style={font=\small},label style={font=\small},xmin=-10.5,xmax=10.5,ymin=1e-5,ymax=1,grid=major,grid style={gray!25},
  legend columns=2,legend style={font=\footnotesize,at={(0.5,-0.32)},anchor=north,draw=none,fill=none,
  /tikz/every even column/.append style={column sep=8pt}},table/col sep=comma,unbounded coords=jump]
\addplot[omDef,only marks,mark=*,mark size=1.3pt] table[x=z,y=freq]{figdata/research/05-stylised-facts-of-returns/hist.csv};
\addplot[omThm,thick,dashed] table[x=z,y=normal]{figdata/research/05-stylised-facts-of-returns/hist.csv};
\legend{Marché américain 1926--2026,loi normale}
\end{semilogyaxis}
\end{tikzpicture}

\omcaption{Fréquence des rendements quotidiens excédentaires du marché américain, par classes d'un demi-écart-type, de juillet 1926 à juillet 2026, face
à la loi normale de même moyenne et même écart-type. La loi normale attribue aux classes au-delà de cinq écarts-types un effectif cumulé
attendu de 0.015~; le siècle en compte 103 jours de ce type. Données~: Kenneth R.~French data library (statistiques dérivées).}\label{fig:rs:stylised-facts-of-returns:hist}
\end{omfigure}

\begin{definition}[Gaussianité par agrégation]\label{def:rs:stylised-facts-of-returns:aggregational}
La \emph{gaussianité par agrégation}\index{gaussianité par agrégation} est la tendance des rendements cumulés sur des horizons plus longs à
ressembler davantage à une loi normale~: leur kurtosis tend vers trois lorsque l'horizon s'allonge, de sorte que la forme de la distribution dépend de
l'horizon.
\end{definition}

La décroissance est lente (la \cref{fig:rs:stylised-facts-of-returns:agg}). Le kurtosis des rendements logarithmiques du siècle vaut 20.0 à un jour,
11.5 à une semaine, 9.5 à un mois et 7.3 à un trimestre. Des tirages indépendants ramèneraient un kurtosis de 20 à $3 + 17/21 = 3.8$ sur un
mois de 21 jours~; les rendements sont non corrélés mais non indépendants, et le regroupement de volatilité maintient les queues mensuelles épaisses.

\begin{omfigure}
\begin{tikzpicture}
\begin{semilogxaxis}[width=10cm,height=5cm,xlabel={horizon (jours de cotation, échelle logarithmique)},ylabel={kurtosis des rendements logarithmiques},
  tick label style={font=\small},label style={font=\small},xmin=0.8,xmax=80,ymin=0,ymax=22,grid=major,grid style={gray!25},
  xtick={1,5,10,21,63},xticklabels={1,5,10,21,63},legend columns=1,legend style={font=\footnotesize,at={(0.5,-0.32)},
  anchor=north,draw=none,fill=none,/tikz/every even column/.append style={column sep=8pt}},table/col sep=comma]
\addplot[omDef,very thick,mark=*] table[x=horizon,y=french]{figdata/research/05-stylised-facts-of-returns/agg.csv};
\addplot[omThm,very thick,dashed,mark=square*] table[x=horizon,y=synth]{figdata/research/05-stylised-facts-of-returns/agg.csv};
\addplot[black!60,thin,sharp plot] coordinates {(0.8,3) (80,3)};
\legend{Marché américain 1926--2026,facteur de marché synthétique,loi normale}
\end{semilogxaxis}
\end{tikzpicture}

\omcaption{Kurtosis des rendements logarithmiques cumulés sur des horizons sans recouvrement. Le kurtosis du marché américain passe de 20.0 à un jour à 9.5 à un
mois et 7.3 à un trimestre~; le facteur de marché synthétique du livre (dix ans, GJR-GARCH avec chocs de Student-$t$) passe de 18.3 à 3.8
à un mois. Données~: Kenneth R.~French data library (statistiques dérivées)~; \texttt{firm.synthmkt}, graine 1.}\label{fig:rs:stylised-facts-of-returns:agg}
\end{omfigure}

\section{Regroupement de volatilité et mémoire longue dans l'ampleur des rendements}

Les rendements sont presque non corrélés~: l'autocorrélation d'ordre un du siècle vaut 0.046. Leur ampleur ne l'est pas. De forts mouvements sont suivis
de forts mouvements (regroupement de volatilité, livre~4, chapitre~18), et la dépendance perdure~: l'autocorrélation des rendements absolus vaut 0.30
à un jour, 0.22 à vingt jours et 0.13 à cent jours (la \cref{fig:rs:stylised-facts-of-returns:acf}). Cont décrit la décroissance comme
à peu près une puissance du décalage, avec un exposant compris entre 0.2 et 0.4, la mémoire longue du livre~4, chapitre~17.

\begin{definition}[Effet Taylor]\label{def:rs:stylised-facts-of-returns:taylor}
L'\emph{effet Taylor}\index{effet Taylor} est le constat selon lequel les autocorrélations des rendements absolus dépassent celles de leurs
puissances $|r|^\theta$ pour $\theta \ne 1$, en particulier celles des rendements au carré (Granger et Ding, 1996, d'après Taylor, 1986).
\end{definition}

Le siècle le montre à tous les décalages~: 0.30 contre 0.26 à un jour, 0.22 contre 0.11 à vingt jours, 0.13 contre 0.04 à cent jours.
Les carrés accordent un poids excessif aux plus fortes variations, et celles-ci forment la part la moins persistante de la volatilité. La
conséquence pratique est modeste mais réelle~: les modèles et les diagnostics de la persistance de la volatilité construits sur les rendements absolus en voient davantage que
ceux construits sur les carrés.

\begin{omfigure}
\begin{tikzpicture}
\begin{semilogxaxis}[width=10.5cm,height=5.2cm,xlabel={décalage (jours de cotation, échelle logarithmique)},ylabel={autocorrélation},
  tick label style={font=\small},label style={font=\small},xmin=1,xmax=250,ymin=-0.05,ymax=0.35,grid=major,grid style={gray!25},
  legend columns=2,legend cell align=left,legend style={font=\footnotesize,at={(0.5,-0.32)},anchor=north,draw=none,fill=none,
  /tikz/every even column/.append style={column sep=8pt}},table/col sep=comma]
\addplot[omDef,very thick,mark=*,mark size=1.2pt] table[x=lag,y=abs]{figdata/research/05-stylised-facts-of-returns/acf.csv};
\addplot[omDef!60,thick,mark=o,mark size=1.2pt] table[x=lag,y=sq]{figdata/research/05-stylised-facts-of-returns/acf.csv};
\addplot[black!60,thick,mark=x,mark size=1.6pt] table[x=lag,y=r]{figdata/research/05-stylised-facts-of-returns/acf.csv};
\addplot[omThm,very thick,dashed] table[x=lag,y=synth_abs]{figdata/research/05-stylised-facts-of-returns/acf.csv};
\legend{$|r|$ (marché américain),$r^2$ (marché américain),$r$ (marché américain),$|r|$ (marché synthétique)}
\end{semilogxaxis}
\end{tikzpicture}

\omcaption{Autocorrélations des rendements quotidiens, des rendements absolus et des rendements au carré du marché américain, 1926--2026, et des
rendements absolus du facteur de marché synthétique. Les rendements sont presque non corrélés~; les rendements absolus restent corrélés au-delà de cent
jours, et davantage que les rendements au carré (l'\omterm{def:rs:stylised-facts-of-returns:taylor}{effet Taylor}). Le regroupement du facteur synthétique s'estompe après environ cinquante jours. Données~:
Kenneth R.~French data library (statistiques dérivées)~; \texttt{firm.synthmkt}, graine 1.}\label{fig:rs:stylised-facts-of-returns:acf}
\end{omfigure}

\section{Effet de levier et asymétrie gains--pertes}

\begin{definition}[Asymétrie gains--pertes]\label{def:rs:stylised-facts-of-returns:gainloss}
L'\emph{asymétrie gains--pertes}\index{asymétrie gains--pertes} est le constat que les cours et les indices boursiers connaissent de fortes baisses mais
pas de hausses aussi fortes~: la distribution des rendements est à asymétrie négative, et davantage aux horizons de plusieurs semaines qu'à ceux de quelques jours.
\end{definition}

L'asymétrie quotidienne du siècle vaut $-0.16$~; sur une semaine elle vaut $-0.76$, et sur un mois $-0.90$. La plus forte baisse quotidienne est celle de $-17.4~\%$
du 19 octobre 1987~; la plus forte hausse quotidienne, $+15.7~\%$ le 15 mars 1933, est intervenue sur le marché le plus volatil du siècle. L'effet de levier
(livre~4, chapitre~18), la corrélation négative entre les rendements et la volatilité ultérieure, est visible dans les mêmes données~: la
corrélation entre le rendement d'aujourd'hui et le carré du rendement de demain vaut $-0.085$, contre $+0.025$ entre le rendement d'aujourd'hui et
le carré du rendement d'hier. Les baisses font monter la volatilité~; la volatilité ne prédit pas le signe des rendements.

\section{Saisonnalité intrajournalière}

\begin{definition}[Saisonnalité intrajournalière, profil de volume intrajournalier]\label{def:rs:stylised-facts-of-returns:intraday}
La \emph{saisonnalité intrajournalière}\index{saisonnalité intrajournalière} est la dépendance régulière de la volatilité, du volume, des fourchettes et
de l'activité de négociation à l'heure de la journée. Le \emph{profil de volume intrajournalier}\index{profil de volume intrajournalier} est la part moyenne du
volume d'une journée échangée dans chaque intervalle de la séance.
\end{definition}

Les marchés d'actions sont les plus actifs, et les plus mobiles, juste après l'ouverture et juste avant la clôture~: Wood, McInish et Ord (1985) ont trouvé
des rendements et des écarts-types anormalement élevés au début et à la fin de la journée, dans des données minute par minute de la NYSE, et Andersen et
Bollerslev (1997) modélisent cette périodicité intrajournalière conjointement avec la persistance de la volatilité d'un jour à l'autre. Le constat importe à quiconque échantillonne en temps d'horloge~: un rendement sur les cinq premières minutes de la journée et un rendement sur
cinq minutes à midi sont des tirages de distributions différentes. Le simulateur intrajournalier du livre construit le profil au chapitre~2~: ce jour-là,
le premier quart d'heure échange 2.6 fois le volume du quart d'heure médian, et le dernier 1.7 fois. Le chapitre ne dispose d'aucune donnée intrajournalière sous licence
pour la mesurer sur des actions réelles~; le profil du simulateur est une hypothèse, énoncée comme telle.

\section{Structure transversale}

Les actions évoluent ensemble, et l'ampleur de ce mouvement dépend de la direction du marché.

\begin{definition}[Corrélation de dépassement, asymétrie des corrélations]\label{def:rs:stylised-facts-of-returns:exceedance}
Une \emph{corrélation de dépassement}\index{corrélation de dépassement} est une corrélation entre des rendements mesurée uniquement les jours où une
variable conditionnante (un rendement de marché, ou les deux rendements eux-mêmes) se situe au-delà d'un seuil. L'\emph{asymétrie des
corrélations}\index{asymétrie des corrélations} est le constat que les corrélations de dépassement sur les marchés en baisse dépassent celles des marchés en hausse
au même seuil.
\end{definition}

Longin et Solnik (2001) ont établi, par des méthodes de valeurs extrêmes, que la corrélation entre les marchés d'actions internationaux augmente sur les marchés
baissiers mais pas sur les marchés haussiers~; Ang et Chen (2002) ont trouvé la même asymétrie entre des portefeuilles d'actions américaines et le marché. La
bibliothèque de données de Kenneth R.~French, à travers ses dix portefeuilles sectoriels, montre ce phénomène sur le siècle (\cref{fig:rs:stylised-facts-of-returns:asym})~: les jours où
le marché a baissé de plus d'un écart-type, leur corrélation moyenne par paires vaut 0.61~; les jours où il a monté autant, 0.56~; l'écart
existe à chaque seuil entre zéro et deux écarts-types. La diversification entre secteurs vaut le moins les jours où
elle est la plus nécessaire.

Le second fait transversal est un facteur dominant. La matrice de corrélation de 200 actions simulées sur un an a une plus grande
valeur propre de 44.5, soit 22~\% du total, contre une borne supérieure de Marchenko--Pastur de 3.58 pour du bruit pur de même forme (livre~4,
chapitre~22)~; quatre valeurs propres dépassent la borne~: le marché et les secteurs les plus forts.

\begin{omfigure}
\begin{tikzpicture}
\begin{axis}[width=10.5cm,height=5.2cm,xlabel={seuil (écarts-types du rendement quotidien du marché)},
  ylabel={corrélation moyenne par paires},tick label style={font=\small},label style={font=\small},xmin=-0.05,xmax=2.05,ymin=0.1,
  ymax=0.7,grid=major,grid style={gray!25},legend columns=2,legend cell align=left,legend style={font=\footnotesize,at={(0.5,-0.32)},anchor=north,
  draw=none,fill=none,/tikz/every even column/.append style={column sep=8pt}},table/col sep=comma]
\addplot[omThm,very thick,mark=*] table[x=threshold,y=down]{figdata/research/05-stylised-facts-of-returns/asym.csv};
\addplot[omDef,very thick,mark=square*] table[x=threshold,y=up]{figdata/research/05-stylised-facts-of-returns/asym.csv};
\addplot[omThm!60,thick,dashed,mark=o] table[x=threshold,y=synth_down]{figdata/research/05-stylised-facts-of-returns/asym.csv};
\addplot[omDef!60,thick,dashed,mark=square] table[x=threshold,y=synth_up]{figdata/research/05-stylised-facts-of-returns/asym.csv};
\legend{10 secteurs américains~: jours de baisse,10 secteurs américains~: jours de hausse,actions synthétiques~: jours de baisse,actions synthétiques~: jours de hausse}
\end{axis}
\end{tikzpicture}

\omcaption{Corrélation moyenne par paires les jours où le marché a baissé (down) ou monté (up) de plus que le seuil. Les dix
portefeuilles sectoriels américains, 1926--2026, sont davantage corrélés les jours de baisse à chaque seuil (0.61 contre 0.56 à un écart-type).
Les actions synthétiques du livre montrent l'ordre inverse~: le simulateur manque ce fait. Données~: bibliothèque de données de Kenneth R.~French
(statistiques dérivées)~; \texttt{firm.synthmkt}, graine 1.}\label{fig:rs:stylised-facts-of-returns:asym}
\end{omfigure}

\section{Bilan du simulateur du livre}

\texttt{firm.synthmkt}, la construction de ce chapitre, engendre la structure transversale que le reste du livre négocie~: un facteur de marché à
volatilité GJR-GARCH et chocs de Student-$t$, des facteurs sectoriels et de style, des actions à volatilité spécifique hétérogène, des sauts liés aux
résultats, et quatre sources de rendement espéré implantées. Le tableau confronte ce simulateur au siècle de données américaines sur les faits de ce chapitre.

\begin{center}
\small
\begin{tabular}{@{}p{5.4cm}cccc@{}}
\toprule
& réel & \multicolumn{3}{c}{synthétique (\texttt{firm.synthmkt})} \\
\cmidrule(lr){2-2}\cmidrule(lr){3-5}
fait (rendements quotidiens) & marché américain & marché & indice EW & action médiane \\
\midrule
kurtosis & 19.1 & 20.5 & 18.0 & 8.7 \\
indice de queue gauche (Hill, 2~\% les plus grandes valeurs) & 2.90 & 2.83 & 2.88 & 3.35 \\
autocorrélation de $|r|$, retard 1 & 0.30 & 0.23 & 0.25 & 0.10 \\
autocorrélation de $|r|$, retard 20 & 0.22 & 0.21 & 0.22 & 0.08 \\
autocorrélation de $|r|$, retard 100 & 0.13 & $-0.02$ & $-0.03$ & $-0.01$ \\
autocorrélation de $r$, retard 1 & 0.05 & $-0.01$ & 0.01 & $-0.03$ \\
corrélation de $r_t$ avec $r^2_{t+1}$ & $-0.085$ & $-0.111$ & $-0.103$ & $-0.024$ \\
kurtosis des rendements logarithmiques mensuels & 9.5 & 3.8 & 3.9 & 3.5 \\
\bottomrule
\end{tabular}
\end{center}

\begin{remark}[Ce que le simulateur peut et ne peut pas servir à faire]\label{rem:rs:stylised-facts-of-returns:scorecard}
Le simulateur reproduit les queues quotidiennes, le regroupement sur un mois et l'effet de levier, ce dont ont besoin les backtests,
les mesures de prédicteurs et les modèles de risque de ce livre. Il manque trois faits~: la mémoire de la volatilité au-delà d'environ cinquante jours (le GARCH
décroît exponentiellement), les queues épaisses des rendements mensuels (il n'a pas de krachs de plus d'une journée), et l'\omterm{def:rs:stylised-facts-of-returns:exceedance}{asymétrie des corrélations}. Les recherches
dont la réponse dépend de ces faits (couverture des queues, risque de drawdown sur plusieurs mois, valeur de la diversification lors d'une chute des marchés) ne doivent pas
y être validées.
\end{remark}

\section{Tutoriel~: évaluer un simulateur}\label{tut:rs:stylised-facts-of-returns:scorecard}

\begin{tutorial}
\textbf{Objectif.} Calculer les \omterm{def:rs:stylised-facts-of-returns:fact}{faits stylisés} d'une série de rendements, les appliquer au siècle de données américaines et au marché synthétique,
et remplir le bilan. \textbf{État final~:} le tableau ci-dessus et les
\cref{fig:rs:stylised-facts-of-returns:agg,fig:rs:stylised-facts-of-returns:acf,fig:rs:stylised-facts-of-returns:asym}.
\begin{tutsteps}
\item \textbf{Les statistiques.} Une fonction calcule chaque ligne du bilan pour une série~; l'estimateur de Hill est celui du livre~4.
\omcode{code/research/05-stylised-facts-of-returns/python/rs_stylised.py}{55}{64}{Le bilan d'une série de rendements.}\label{lst:rs:stylised-facts-of-returns:facts}
\item \textbf{Le facteur de marché.} GJR-GARCH(1,1)~: les chocs négatifs font monter la variance du lendemain davantage que les chocs positifs~; les
chocs de Student-$t$ donnent les queues que le seul regroupement n'explique pas.
\omcode{code/firm/synthmkt/firm_synthmkt.py}{156}{167}{Le facteur de marché synthétique.}\label{lst:rs:stylised-facts-of-returns:garch}
\item \textbf{Exécuter} \texttt{synth\_series()} pour les lignes synthétiques, \texttt{french\_summary()} pour la ligne réelle, et
\texttt{fig\_stylised.py}. Les statistiques réelles proviennent de \texttt{rs\_fetch\_french.py}, qui ne conserve que les statistiques dérivées des
données de la bibliothèque.
\end{tutsteps}
\textbf{Pour aller plus loin.} Porter \texttt{garch\_beta} à 0.96 et \texttt{garch\_alpha} à 0.01 (même persistance, décroissance plus lente)
et observer le déplacement de l'autocorrélation au retard 100~; rendre les bêtas plus grands les jours de baisse et chercher l'\omterm{def:rs:stylised-facts-of-returns:exceedance}{asymétrie des corrélations}, et ce qu'elle fait
au rendement moyen.
\end{tutorial}

\section{Construction~: le marché synthétique d'actions}\label{bld:rs:stylised-facts-of-returns:synthmkt}

\begin{build}
\bfield{Rôle} Le panier de référence du livre, avec la vérité attachée. Les chapitres 6 à 15 y mesurent les prédicteurs, la partie IV y fait les backtests,
la partie V y construit les modèles de risque et les portefeuilles, et les livres 8 et 9 y mènent leurs tutoriels sur les actions.
\bfield{Interface} \texttt{MarketConfig} (taille, longueur, graine~; paramètres des facteurs de marché, sectoriels et de style~; volatilité spécifique~; les
quatre forces d'alpha injectées~; règles de radiation, de fusion, de division, de dividende et de dépôt de documents)~; \texttt{simulate(cfg)} renvoyant un \texttt{Panel}~:
\texttt{ret}, \texttt{price}, \texttt{volume}, \texttt{shares}, \texttt{listed} quotidiens~; \texttt{industry}, \texttt{beta},
\texttt{style\_x}, \texttt{spec\_vol}~; les trajectoires des facteurs et la variance conditionnelle du marché~; \texttt{alpha} (le rendement espéré injecté
pour le lendemain, par composante)~; les radiations, les divisions, les dividendes, les annonces de résultats avec leurs surprises véritables, les fondamentaux en temps réel
avec leurs \omterm{def:rs:point-in-time-data-and-the-biases:vintage}{retraitements}~; un \texttt{secmaster}.
\bfield{Règles} Déterministe pour une graine donnée. Le rendement espéré injecté pour le jour $t+1$ est enregistré au jour $t$, de sorte qu'un signal connu à
la clôture de $t$ puisse être évalué par rapport à lui. Un titre radié est remplacé le lendemain. Les fondamentaux sont publiés après la fin du trimestre,
et les \omterm{def:rs:point-in-time-data-and-the-biases:vintage}{retraitements} arrivent plus tard encore.
\bfield{Tests d'acceptation} \texttt{code/firm/synthmkt/tests/}~: déterminisme et tenue des comptes (nombre de titres cotés, \omterm{def:rs:universe-symbology-and-corporate-actions:delisting}{rendements de radiation}, aucun rendement
après une radiation)~; accord avec le \omterm{def:rs:universe-symbology-and-corporate-actions:master}{référentiel des titres} (security master)~; fondamentaux en temps réel via \texttt{firm.pit}~; sur trois graines, pour le marché
kurtosis supérieure à 5, regroupement de volatilité positif à vingt jours, corrélation de levier négative, et les
coefficients d'information positifs pour le retour à la moyenne et le momentum~; une première valeur propre dominante.
\bfield{Pour aller plus loin} Bêtas à la baisse avec une dérive compensatrice (\omterm{def:rs:stylised-facts-of-returns:exceedance}{asymétrie des corrélations})~; une composante de volatilité
qui varie lentement (mémoire longue)~; krachs sur plusieurs jours.
\end{build}

\begin{omsources}
\item R.~Cont, ``Empirical properties of asset returns: stylized facts and statistical issues'', \emph{Quantitative Finance} 1(2),
2001.
\item P.~Gopikrishnan, M.~Meyer, L.~A.~N.~Amaral and H.~E.~Stanley, ``Inverse cubic law for the distribution of stock price
variations'', \emph{European Physical Journal B} 3, 1998.
\item C.~W.~J.~Granger and Z.~Ding, ``Varieties of long memory models'', \emph{Journal of Econometrics} 73, 1996; S.~J.~Taylor,
\emph{Modelling Financial Time Series}, Wiley, 1986.
\item R.~A.~Wood, T.~H.~McInish and J.~K.~Ord, ``An investigation of transactions data for NYSE stocks'', \emph{Journal of Finance}
40(3), 1985; T.~G.~Andersen and T.~Bollerslev, ``Intraday periodicity and volatility persistence in financial markets'',
\emph{Journal of Empirical Finance} 4, 1997.
\item F.~Longin and B.~Solnik, ``Extreme correlation of international equity markets'', \emph{Journal of Finance} 56(2), 2001;
A.~Ang and J.~Chen, ``Asymmetric correlations of equity portfolios'', \emph{Journal of Financial Economics} 63(3), 2002.
\item B.~Mandelbrot, ``The variation of certain speculative prices'', \emph{Journal of Business} 36(4), 1963.
\item Kenneth R.~French, Data Library: Fama/French 3 factors and 10 industry portfolios, daily.
\item Federal Reserve History, ``Stock Market Crash of 1987''.
\end{omsources}

\section{Exercices}

\begin{exercise}[$\star$]\label{exo:rs:stylised-facts-of-returns:1}
Une loi normale d'écart type quotidien du siècle de 1.08~\%~: quelle est la probabilité d'un jour sous $-17.4~\%$, et combien de tels jours
attend-on sur 26\,296~?
\end{exercise}

\begin{exercise}[$\star$]\label{exo:rs:stylised-facts-of-returns:2}
Si les rendements quotidiens étaient indépendants de kurtosis 20, quelle serait la kurtosis de leur somme sur 5, 21 et 63 jours~? Comparer avec les
11.5, 9.5 et 7.3 du siècle.
\end{exercise}

\begin{exercise}[$\star$]\label{exo:rs:stylised-facts-of-returns:3}
Une queue en loi de puissance d'indice 3~: si un jour sur mille chute de plus de 4~\%, à quelle fréquence un jour chute-t-il de plus de 8~\%~? Et avec un
indice de 4~?
\end{exercise}

\begin{exercise}[$\star\star$]\label{exo:rs:stylised-facts-of-returns:4}
Montrer que pour une loi de Student-$t$ à $\nu > 4$ degrés de liberté, normalisée à variance unitaire, la kurtosis vaut $3 + 6/(\nu - 4)$. Quel $\nu$
correspond à une kurtosis de 19~?
\end{exercise}

\begin{exercise}[$\star\star$]\label{exo:rs:stylised-facts-of-returns:5}
Les dix secteurs ont une corrélation moyenne par paires de 0.61 les jours de baisse du marché et de 0.56 les jours de hausse (à un
écart type). Avec des poids égaux et des volatilités égales $\sigma$, quelle est la volatilité du portefeuille à pondération égale relativement à
$\sigma$ dans chacun des deux cas~?
\end{exercise}

\begin{exercise}[$\star\star$]\label{exo:rs:stylised-facts-of-returns:6}
Le facteur de marché GJR-GARCH a $\alpha = 0.02$, $\gamma = 0.10$, $\beta = 0.92$. Quelle est sa persistance, et la demi-vie d'un choc de
variance~? Pourquoi son autocorrélation des rendements absolus s'éteint-elle après une cinquantaine de jours~?
\end{exercise}

\begin{exercise}[$\star\star\star$]\label{exo:rs:stylised-facts-of-returns:7}
\emph{Programmation.} Simuler le marché synthétique avec \texttt{garch\_alpha=0.01}, \texttt{garch\_gamma=0.05},
\texttt{garch\_beta=0.965} (la même persistance) et comparer les autocorrélations des rendements absolus du marché aux retards 20 et 100
avec celles du réglage par défaut.
\end{exercise}

\begin{exercise}[$\star\star\star$]\label{exo:rs:stylised-facts-of-returns:8}
\emph{Trouver la faille.} «~Nous avons validé notre stratégie de couverture de queue sur un marché simulé calibré sur la kurtosis quotidienne et le regroupement de
volatilité du S\&P~500. Elle se rentabilise sur toute fenêtre de dix ans.~»
\end{exercise}

\section{Problème~: le simulateur est-il réaliste~?}

\begin{problem}[{Problème du week-end --- bilan du marché synthétique du livre face à un siècle de données}]\label{pb:rs:stylised-facts-of-returns:1}
Le siècle de rendements en excès quotidiens du marché américain (la Data Library de French, statistiques dérivées) et \texttt{firm.synthmkt} avec sa
configuration par défaut (1\,000 titres, dix ans, graine 1).

\medskip\noindent\textbf{Partie I --- Queues.}
\begin{enumerate}
\item Quelles sont les kurtoses quotidiennes du marché américain et du facteur de marché synthétique~?
\item Quels indices de queue gauche l'estimateur de Hill donne-t-il à chacun sur ses 2~\% de pertes les plus fortes~?
\item Combien de jours du siècle se situent au-delà de cinq écarts types, et combien une loi normale en prévoirait-elle~?
\item À combien d'écarts types des cinq années précédentes se situait le 19 octobre 1987~?
\item Quelle kurtosis a le titre synthétique médian, et pourquoi est-elle inférieure à celle de l'indice~?
\end{enumerate}

\medskip\noindent\textbf{Partie II --- Dépendance.}
\begin{enumerate}[resume]
\item Quelles sont les autocorrélations des rendements absolus aux retards 1, 20 et 100, réels et synthétiques~?
\item Où le marché synthétique échoue-t-il, et pourquoi~?
\item Le marché synthétique présente-t-il l'effet de levier~? Comparer les corrélations de $r_t$ avec $r^2_{t+1}$.
\item Quelles sont les kurtoses des rendements logarithmiques mensuels, et qu'est-ce qui explique l'écart~?
\item Le siècle présente-t-il l'\omterm{def:rs:stylised-facts-of-returns:taylor}{effet Taylor} au retard 20~?
\end{enumerate}

\medskip\noindent\textbf{Partie III --- La section transversale.}
\begin{enumerate}[resume]
\item Quelles sont les corrélations moyennes des dix secteurs les jours de baisse et de hausse à un écart type~?
\item Que donne le marché synthétique, et qu'est-ce que cela signifie~?
\item Quelle est la plus grande valeur propre de la corrélation de 200 titres synthétiques sur une année, sa part, et le bord de Marchenko--Pastur~?
\item Combien de valeurs propres se situent au-dessus du bord, et que sont-elles~?
\end{enumerate}

\medskip\noindent\textbf{Partie IV --- Verdict.}
\begin{enumerate}[resume]
\item Quels travaux de cet ouvrage peuvent utiliser le simulateur sans réserve~?
\item Quels travaux ne doivent pas être validés à l'aide de celui-ci~?
\item Quel changement unique l'améliorerait le plus pour une étude du risque de queue~?
\item Énoncer le \emph{résultat nommé}~: les trois lignes de la grille de contrôle que le simulateur reproduit le mieux et les trois qu'il manque.
\item Pourquoi un simulateur utilisé pour la recherche a-t-il besoin d'une grille de contrôle~?
\item En une phrase~: à quoi sert un \omterm{def:rs:stylised-facts-of-returns:fact}{fait stylisé}~?
\end{enumerate}
\end{problem}

\section{Questions d'entretien}

\begin{interviewq}[$\star$ \iqroles{chercheur, trader}]\label{iq:rs:stylised-facts-of-returns:1}
Citez quatre \omterm{def:rs:stylised-facts-of-returns:fact}{faits stylisés} des rendements d'actions.
\end{interviewq}

\begin{interviewq}[$\star\star$ \iqroles{chercheur, risque}]\label{iq:rs:stylised-facts-of-returns:2}
Les rendements quotidiens sont presque non corrélés. Sont-ils pour autant indépendants~? Que vérifieriez-vous~?
\end{interviewq}

\begin{interviewq}[$\star\star$ \iqroles{risque, chercheur}]\label{iq:rs:stylised-facts-of-returns:3}
Pourquoi la diversification sectorielle déçoit-elle lors d'un krach~?
\end{interviewq}

\begin{interviewq}[$\star\star$ \iqroles{chercheur, ingénieur ML}]\label{iq:rs:stylised-facts-of-returns:4}
Vous voulez un simulateur pour tester une stratégie de trading. Comment décidez-vous s'il est suffisamment réaliste~?
\end{interviewq}

\begin{interviewq}[$\star\star\star$ \iqroles{chercheur}]\label{iq:rs:stylised-facts-of-returns:5}
Si les rendements quotidiens ont une kurtosis de 20 et sont indépendants, quelle est la kurtosis des rendements mensuels~? Les données donnent environ 9. Qu'est-ce que cela
vous apprend~?
\end{interviewq}

\begin{interviewq}[$\star\star$ \iqroles{trader, chercheur}]\label{iq:rs:stylised-facts-of-returns:6}
Qu'est-ce que l'effet de levier, et comment le mesureriez-vous~?
\end{interviewq}
